% !PRINTED 543
\weberpage{543}
%<<CONT
$\pi_i$ theilbar, durch andere nicht theilbar sein soll, und dieser
Forderung kann nach \S.\,143, 3. immer genügt werden. Dann können wir
nach 2. die ganze Zahl $\tau$ so bestimmen, dass
\[
\alpha\xi + \beta_1\tau = 1
\]
wird, und wenn wir $\tau\eta_1 = \eta$, $\tau\zeta_1 = \zeta, \ldots$
setzen, so erhalten wir den Satz:

3. \gesperrt{Sind $\alpha, \beta, \gamma, \ldots$ Zahlen in
$\mathfrak{o}$ ohne gemeinschaftlichen Theiler, so lassen sich andere
Zahlen $\xi, \eta, \zeta, \ldots$ in $\mathfrak{o}$ so bestimmen, dass}
\weq{4}{\alpha\xi + \beta\eta + \gamma\zeta + \cdots = 1}
\gesperrt{wird.}

Daraus lässt sich weiter auf folgenden Satz schliessen:

4. \gesperrt{Sind $\alpha, \beta, \gamma, \ldots$ beliebige Zahlen in
$\mathfrak{o}$, $\mu$ eine durch den grössten gemeinschaftlichen Theiler
aller dieser Zahlen theilbare Zahl in $\mathfrak{o}$, so kann man die
Zahlen $\xi, \eta, \zeta, \ldots$ in $\mathfrak{o}$ so bestimmen, dass}
\weq{5}{\mu = \alpha\xi + \beta\eta + \gamma\zeta + \cdots}
\gesperrt{wird.}

Wenn wir nämlich den grössten gemeinschaftlichen Theiler der Zahlen
$\alpha, \beta, \gamma, \ldots$ nach \S.\,139 in der Form
\[
\delta = \alpha u + \beta v + \gamma w \ldots
\]
annehmen, worin $u, v, w, \ldots$ Variable sind, so giebt es nach
Voraussetzung ein ganzes Functional $\omega$, so dass $\mu = \delta
\omega$ ist. Das Functional $\omega$ stellen wir als Quotienten zweier
ganzer Functionen $\varphi : \varepsilon$ dar, von denen $\varepsilon$
eine Einheit, und folglich $\varphi$ eine Function mit ganzzahligen
Co\"efficienten ist, und erhalten so
\weq{6}{\varepsilon\mu = (\alpha u + \beta v + \gamma w + \cdots)\varphi.}
Ordnet man beide Seiten dieser Gleichung nach den darin vorkommenden
Variablen, so erhält man, wenn $\varepsilon_1, \varepsilon_2, \ldots$
die Co\"efficienten in $\varepsilon$ sind, ein System von Gleichungen
von folgender Form:
\weq{7}{\begin{aligned}
\varepsilon_1\mu &= \alpha\xi_1 + \beta\eta_1 + \gamma\zeta_1 + \cdots \\
\varepsilon_2\mu &= \alpha\xi_2 + \beta\eta_2 + \gamma\zeta_2 + \cdots \\
& \text{.\,.\,.\,.\,.\,.\,.\,.\,.\,.\,.\,.\,.\,.\,.\,.\,.\,.\,.\,.\,.\,.\,.,}
\end{aligned}}
worin die $\xi_i, \eta_i, \zeta_i, \ldots$ Zahlen in $\mathfrak{o}$ sind.
Nun haben aber die Zahlen $\varepsilon_1, \varepsilon_2, \ldots$ als
Co\"efficienten einer Einheit, keinen gemein
%CONT>>
